% 
%
% (c) 2025 Autor, ETH Zürich
%
% !TEX root = linalg_I-II_summary.tex
% !TEX encoding = UTF-8
%

\section{Orthogonal \& Unitary maps / Isometries}
\begin{definition}{Linear Isomerty}{isom}
	Let $V, W$ be inner product spaces over $K = \mathbb{R}$ or $\mathbb{C}$. A map $T: V \to W$ is called a \textbf{linear isometry} if $T$ is linear and
	\[
		\forall \, v_1, v_2 \in V : \; \langle Tv_1, Tv_2 \rangle_W = \langle v_1, v_2 \rangle_V,
	\]
	i.e., $T$ preserves the inner product.
\end{definition}

\begin{remark}
	If $T:V \to W$ is a linear isometry, then:
	\begin{enumerate}
		\item $\forall \, v \in V : \; \Vert Tv \Vert = \Vert v \Vert$.
		\item $\forall\, p,q \in V : \; d(Tp, Tq) = d(p, q)$, where $d(p, q) := \Vert p - q \Vert$.
		\item $T$ is injective. Indeed, if $v \in \ker T$, then 
		\[
			0 = \Vert Tv \Vert = \Vert v \Vert \quad \Longrightarrow \quad v = 0.
		\]
		So, $\ker T = \{0\}$.
	\end{enumerate}
\end{remark}

From now on we'll concentrate on the case $W = V$, i.e., $T \in \operatorname{End}(V)$. If such a map $T$ is an isometry we also call it \textbf{orthogonal endomorphism} if $K = \mathbb{R}$, or \textbf{unitary endomorphism} if $K = \mathbb{C}$.

\begin{theorem}{Equivalent Statements to Isometry}{equiv_isom}
	Let $V$ be a finite dimensional inner product space over $K = \mathbb{R}$ or $\mathbb{C}$, and $T \in \operatorname{End}(V)$. The following are equivalent:
	\begin{enumerate}
		\item $T$ is an isometry.
		\item $\forall\, v \in V : \; \Vert Tv \Vert = \Vert v \Vert$.
		\item For every orthonormal system $e_1, \hdots , e_k$ of $V$, the list of vectors $Te_1, \hdots , Te_k$ is also an orthonormal system.
		\item There exists an orthonormal basis $e_1, \hdots , e_n$ of $V$ s.t. $Te_1, \hdots , Te_n$ is also an orthonormal basis.
		\item $T\,T^* = id_V$.
		\item $T^*T = id_V$.
		\item $T^* = T^{-1}$.
		\item $T^*$ is an isometry.
	\end{enumerate}
\end{theorem}

Back to matrices.

\begin{proposition}{Equivalent Condition for a Orthogonal/Unitary Map}{equiv_cond_orth_map}
	Let $V$ be an n-dimensional inner product space over $K = \mathbb{R}$ (respectively $\mathbb{C}$). Let $\mathcal{B}$ be an orthonormal basis for $V$. Let $T \in \operatorname{End}(V)$. Then $T$ is orthogonal (respectively unitary) if and only if $[T]^{\mathcal{B}}_{\mathcal{B}} \in O(n)$ (respectively $[T]^{\mathcal{B}}_{\mathcal{B}} \in U(n)$).
\end{proposition}

Note that $\forall \, A \in O(n)$ we have $\det A = \pm 1$. Indeed,
\[
	1 = \det I = \det (A \, A^T) = \det A \cdot \det A.
\]
Similarly, $\forall \, A \in U(n)$, $\vert \underbrace{\det A}_{\displaystyle \in \mathbb{C}} \vert = 1$. Indeed,
\[
	1 = \det (A \, A^*) = \det A \cdot \det \overline{A}^T = \det A \cdot \det \overline{A} = \det A \cdot \overline{\det A} = \vert \det A \vert ^2.
\]

\begin{definition}{Special Orthogonal/Unitary Group}{special_orth_grp}
	The set
	\[
		SO(n) := \{A \in O(n) \;|\; \det A = 1\} \subseteq O(n) \subseteq GL_n(\mathbb{R})
	\]
	is called the \textbf{special orthogonal group}.
	
	The set
	\[
		SU(n) := \{A \in U(n) \;|\; \det A = 1\} \subseteq U(n) \subseteq GL_n(\mathbb{C})
	\]
	is called the \textbf{special unitary group}.
\end{definition}

\begin{proposition}{Group Structure of $O(n)$/$SO(n)$/$U(n)$/$SU(n)$}{grp_struc_son}
	Let $G = O(n), SO(n), U(n), SU(n)$. Then:
	\begin{enumerate}
		\item $I_n \in G$.
		\item $\forall \, A,B \in G : \; A\cdot B \in G$.
		\item $\forall \, A \in G :\; A^{-1} \in G$.
	\end{enumerate}
\end{proposition}

\subsection{Classification of the Elements of $O(2), O(3)$ etc.}
\begin{lemma}{Eigenvalue of an Orthogonal Matrix}{eigval_orth_mat}
	Let $A \in O(n)$. If $\lambda$ is an eigenvalue of $A$, then $\lambda = \pm 1$.
\end{lemma}

\begin{proposition}{Classification of the Elements of $O(2)$}{rot_mat}
	Let $A \in O(2)$.
	\begin{enumerate}
		\item If $\det A = 1$, then $\exists \, \theta \in [0, 2\pi)$ s.t.
		\[
			A = \begin{pmatrix}
				\cos \theta & -\sin \theta\\
				\sin \theta & \cos \theta
			\end{pmatrix} := R_{\theta}.
		\]
		The map $T_A:\mathbb{R}^2 \to \mathbb{R}^2$ is an anti-clockwise rotation by the angle $\theta$.
		\item If $\det A = -1$, then there exists a basis $\mathcal{B} = (v_1, v_2)$ for $\mathbb{R}^2$ s.t. $T_A: \mathbb{R}^2 \to \mathbb{R}^2$ has the following representation in the basis $\mathcal{B}$
		\begin{equation}
			\label{eq:mirror_mat}
			[T_A]^{\mathcal{B}}_{\mathcal{B}} = \begin{pmatrix}
				1 & 0\\
				0 & -1
			\end{pmatrix}.
		\end{equation}
		In other words $T_A$ performs a reflection w.r.t. the line $l = \operatorname{Sp}(v_1)$. Moreover, $\exists \, \theta \in [0, 2\pi)$ s.t. 
		\[
			A = \begin{pmatrix}
				\cos \theta & \sin \theta\\
				\sin \theta & - \cos \theta
			\end{pmatrix} 
			= R_{\theta} \cdot \begin{pmatrix}
				1 & 0 \\
				0 & -1
			\end{pmatrix}.
		\]
		The basis $\mathcal{B}$ for which Equation \eqref{eq:mirror_mat} holds, is given by
		\[
			v_1 = \begin{pmatrix}
				\cos \tfrac{\theta}{2}\\
				\sin \tfrac{\theta}{2}
			\end{pmatrix}, \, 
			v_2 = \begin{pmatrix}
				\cos(\tfrac{\theta}{2} + \tfrac{\pi}{2})\\
				\sin (\tfrac{\theta}{2} + \tfrac{\pi}{2})
			\end{pmatrix}.
		\]
	\end{enumerate}
\end{proposition}

\textbf{Summary:}
\begin{align*}
	SO(2) &= \text{ rotations}\\
	O(2) \setminus SO(2) &= \text{ reflections}.
\end{align*}

\begin{corollary}{}{cor_rot_mat}
	Let $V$ be a 2-dimensional Euclidean space and $T:V \to V$ a linear isometry. Then $\det T = \pm 1$. Moreover, there exists an orthonormal basis $\mathcal{B}$ for $V$ s.t.
	\begin{enumerate}
		\item If $\det T = 1$, then $[T]^{\mathcal{B}}_{\mathcal{B}} = R_{\theta}$.
		\item If $\det T = -1$, then
		\[
			[T]^{\mathcal{B}}_{\mathcal{B}} = \begin{pmatrix}
				1 & 0 \\
				0 & -1
			\end{pmatrix}.
		\]
	\end{enumerate}
\end{corollary}

\begin{proposition}{\textcolor{red}{Classification of the Elements of $SO(3)$}}{clas_so3}
	Let $A \in SO(3)$. Then 1 is an eigenvalue of $A$. Moreover, there exists an orthonormal basis $\mathcal{B} = (v_1, v_2, v_3)$ and $\theta \in [0, 2\pi)$ s.t.
	\[
		[T_A]^{\mathcal{B}}_{\mathcal{B}} = 
		\left(
		\begin{array}{c|cc}
			1 & 0 & 0 \\
			\hline
			0 & \multicolumn{2}{c}{\smash{\lower2ex\hbox{$R_\theta$}}} \\
			0 & & 
		\end{array}
		\right).
	\]
	So, $T_A$ is a rotation with angle $\theta$ around the axis $\operatorname{Sp}(v_1)$. Moreover,
	\begin{enumerate}
		\item[(a)] If $-1$ is an eigenvalue of $A$, then $\theta = \pi$, hence
		\[
			[T_A]^{\mathcal{B}}_{\mathcal{B}} = \begin{pmatrix}
				1 & 0& 0\\
				0 & -1 & 0\\
				0 & 0 & -1
			\end{pmatrix}.
		\]
		\item[(b)] If -1 is \textbf{not} an eigenvalue of $A$ ($\Longleftrightarrow$ 1 is the only eigenvalue), then $\theta \neq \pi$.
	\end{enumerate}
\end{proposition}

\begin{proof}
	The characteristic polynomial $p_A(x)$ of $A$ is a polynomial of degree 3 and with real coeffiecients. By the intermediate value theorem (Analysis I) there exists a real zero $\lambda \in \mathbb{R}$(,i.e., $p_A(\lambda) = 0$). Therefore, $A$ has at least one real eigenvalue. By Lemma \ref{lem*eigval_orth_mat}, $\lambda = \pm 1$. Let $v_1 \in \mathbb{R}^3 \setminus \{0\}$ be an eigenvector of $\lambda$. We can normalize $v_1$ and still have an eigenvector for $\lambda$. So, let's assume $\Vert v \Vert = 1$. Consider the subspace
	\[
		L := v_1^{\perp} = \{v \in \mathbb{R}^3 \; |\; \langle v, v_1 \rangle=0 \}.
	\]
	Since $A$ is orthogonal, we have that
	\[
		\forall \, v \in L :\; Av \in L,
	\]
	i.e., $T_A(L) \subseteq L$, i.e., $L$ is invariant under $T_A$. Indeed, if $v \in L$, then since $A \in O(3)$
	\begin{align*}
		0 &= \langle v, v_1 \rangle = \langle Av, Av_1\rangle = \langle Av, \pm v_1\rangle = \pm \langle Av, v_1\rangle\\
		\Longrightarrow \quad Av &\in L.
	\end{align*}
	We also have $\mathbb{R}^3 = \operatorname{Sp}(v_1) \oplus L$, hence $\dim L = 2$.
	
	Now $T_A|_{L}:L \to L$ is a linear isometry. Let $\mathcal{L} = (w_1, w_2)$ be an orthonormal basis for $L$. By Proposition \ref{prop*equiv_cond_orth_map}, we have that
	\[
		Q := [T_A|_{L}]^{\mathcal{L}}_{\mathcal{L}} \in O(2).
	\]
	
	\underline{Case (a):} $\lambda = 1$. Take $\mathcal{B} = (v_1, w_1, w_2)$. Then
	\[
		[T_A]^{\mathcal{B}}_{\mathcal{B}} = 
		\left(
		\begin{array}{c|cc}
			1 & 0 & 0 \\
			\hline
			0 & \multicolumn{2}{c}{\smash{\lower2ex\hbox{$Q$}}} \\
			0 & & 
		\end{array}
		\right).
	\]
	Now, since $A \in SO(3)$, we have
	\begin{align*}
		1 &= \det A = \det(T_A) = 1 \cdot \det Q\\
		\Longrightarrow \quad \det Q = 1.
	\end{align*}
	So, $Q \in SO(2)$. By Proposition \ref{prop*clas_so3}, $\exists \, \theta \in [0, 2\pi)$ s.t. $Q = R_{\theta}$. Therefore,
	\[
		[T_A]^{\mathcal{B}}_{\mathcal{B}} = 
		\left(
		\begin{array}{c|cc}
			1 & 0 & 0 \\
			\hline
			0 & \multicolumn{2}{c}{\smash{\lower2ex\hbox{$R_\theta$}}} \\
			0 & & 
		\end{array}
		\right).
	\]
	
	\underline{Case (b):} $\lambda = -1$. By the same argument as for case (a), we get that
	\begin{align*}
		1 &= \det A = (-1)\cdot \det (T_A|_{L})\\
		\Longrightarrow \quad \det(T_A|_{L}) = -1.
	\end{align*}
	By Corollary \ref{cor*cor_rot_mat}, there exists an orthonormal basis $\mathcal{L} = (w_1, w_2)$ for $L$ s.t.
	\[
		[T_A|_{L}]^{\mathcal{L}}_{\mathcal{L}} = \begin{pmatrix}
			1 & 0\\
			0 & -1
		\end{pmatrix}.
	\]
	Take $\mathcal{B}' = (v_1, w_1, w_2)$. Then
	\[
		[T_A]^{\mathcal{B}'}_{\mathcal{B}'} = \begin{pmatrix}
			-1 & 0 & 0\\
			0 & 1 & 0\\
			0 & 0 & -1
		\end{pmatrix}.
	\]
	And for $\mathcal{B} = (w_1, v_1, w_2)$ we have
	\[
		[T_A]^{\mathcal{B}}_{\mathcal{B}} = 
		\begin{pmatrix}
			1 & 0 & 0\\
			0 & -1 & 0\\
			0 & 0 & -1
		\end{pmatrix}
		=
		\left(
		\begin{array}{c|cc}
			1 & 0 & 0 \\
			\hline
			0 & \multicolumn{2}{c}{\smash{\lower2ex\hbox{$R_\pi$}}} \\
			0 & & 
		\end{array}
		\right).
	\]
	Note that in case (b), 1 is an eigenvalue of $A$ as well. \qedhere
\end{proof}

\begin{proposition}{Classification of $O(3) \setminus SO(3)$}{clas_o3minusso3}
	Let $A \in O(3)\setminus SO(3)$. Then there exists an orthonormal basis $\mathcal{B} = (v_1, v_2, v_3)$ for $\mathbb{R}^3$ s.t.
	\[
		[T_A]^{\mathcal{B}}_{\mathcal{B}}
		=
		\left(
		\begin{array}{c|cc}
			-1 & 0 & 0 \\
			\hline
			0 & \multicolumn{2}{c}{\smash{\lower2ex\hbox{$R_\theta$}}} \\
			0 & & 
		\end{array}
		\right)
	\]
	for some $\theta \in [0, 2\pi)$. (So $T_A$ does a rations by $\theta$ around $\operatorname{Sp}(v_1)$, followed by a reflection w.r.t. $\operatorname{Sp}(v_2, v_3)$.)
\end{proposition}














